% ============================================================
%  Semestrální projekt
%  Fakulta elektrotechnická (FEL, F3), ČVUT v Praze
%  Katedra telekomunikační techniky (K13137)
% ============================================================
%  Překlad: latexmk -pdf soubor.tex

% Volby třídy: twoside = oboustranný tisk; dále lze: oneside, draft
\documentclass[twoside]{ctuthesis}
\usepackage{multirow}
\usepackage{array}
\usepackage{lua-visual-debug}
\usepackage{tabularx}

\ctusetup{
	% --- jazyk práce ---------------------------------------
	mainlanguage   = czech,          % hlavní jazyk textu práce
	otherlanguages = {english},      % další použité jazyky
	titlelanguage  = czech,          % jazyk titulní strany
	% --- typ a zařazení práce ------------------------------
	doctype = B,                     % B = bakalářská, M = diplomová, D = disertační
	faculty = F3,                    % F3 = FEL
	department-czech   = {Katedra telekomunikační techniky},
	department-english = {Department of Telecommunication Engineering},
	% --- názvy práce (PŘESNĚ podle zadání z KOS!) ----------
	title-czech   = {Možnosti spolupráce telefonního systému Asterisk s digitálním radiovým protokolem DMR},
	title-english = {Options for Integrating the Asterisk Telephony System with the DMR Digital Radio Protocol},
	% subtitle-czech   = {Podtitul},        % nepovinné
	% subtitle-english = {Subtitle},        % nepovinné
	% --- autor a vedoucí -----------------------------------
	author             = {Adam Barborík},
	supervisor         = {Ing. Pavel Troller, Csc.},
	supervisor-address = {Katedra telekomunikační techniky, \\ FEL ČVUT v~Praze, \\ Technická~2, 166~27 Praha~6},
	% supervisor-specialist = {Jméno konzultanta},  % nepovinné
	% --- studijní program / specializace -------------------
	fieldofstudy-czech   = {Softwarové inženýrství a technologie},
	fieldofstudy-english = {Software Engineering and Technology},
	% subfieldofstudy-czech   = {Specializace},     % je-li specializace
	% subfieldofstudy-english = {Specialization},
	% --- klíčová slova -------------------------------------
	keywords-czech   = {IP telefonie, Asterisk, DMR, digitální rádio, BrandMeister, VoIP},
	keywords-english = {IP telephony, Asterisk, DMR, digital radio, BrandMeister, VoIP},
	% --- datum odevzdání -----------------------------------
	day = 5, month = 1, year = 2027,
	% --- zadání z KOS (dvoustránkové PDF) ------------------
	% specification-file = {zadani.pdf},   % odkomentuj, až budeš mít PDF zadání z KOS
}

\ctuprocess

% ------------------------------------------------------------
% Definice "přední" dvousloupcové strany. PRÁVĚ ZDE vznikají ty
% dva sloupce s barevnou (fakultní) dělicí linkou – týká se jen
% poděkování, prohlášení, abstraktů, obsahu a seznamu obrázků.
% Samotné kapitoly jsou jednosloupcové.
% ------------------------------------------------------------
\ctutemplateset{maketitle twocolumn default}{
	\begin{twocolumnfrontmatterpage}
		\ctutemplate{twocolumn.thanks}
		\ctutemplate{twocolumn.declaration}
		\ctutemplate{twocolumn.abstract.in.titlelanguage}
		\ctutemplate{twocolumn.abstract.in.secondlanguage}
		\ctutemplate{twocolumn.tableofcontents}
		\ctutemplate{twocolumn.listoffigures}
		% \ctutemplate{twocolumn.listoftables}   % chceš-li i seznam tabulek
	\end{twocolumnfrontmatterpage}
}

% --- Abstrakt česky -------------------------------------------
\begin{abstract-czech}
Sem napiš český abstrakt (jeden odstavec, obvykle 150–300 slov):
téma a kontext, cíl práce, metodika a data, hlavní výsledky.
\end{abstract-czech}

% --- Abstrakt anglicky ----------------------------------------
\begin{abstract-english}
Write the English abstract here (one paragraph).
\end{abstract-english}

% --- Poděkování -----------------------------------------------
\begin{thanks}
Děkuji vedoucímu práce za \ldots
\end{thanks}

% --- Prohlášení -----------------------------------------------
\begin{declaration}
Prohlašuji, že jsem předloženou práci vypracoval samostatně a~že jsem
uvedl veškerou použitou literaturu.

V~Praze, \ctufield{day}.~\monthinlanguage{title}~\ctufield{year}
\end{declaration}

% ==============================================================
\begin{document}

\maketitle
%\twocolumn

\chapter{Vývoj a současný stav digitální rádiové komunikace}

Radioamatérská služba patří mezi nejstarší oblasti využití rádiového spektra. Od vynálezu bezdrátového telegrafu je prostorem pro experimentování
a technické vzdělávání, které ovšem podléhá řadě omezení a regulací, viz kapitola \hyperref[chap:regule]{\ref*{chap:regule} \nameref*{chap:regule}}.
Kromě amatérů je značná část spektra vyhrazená pro komerční nebo profesionální použití. Jak v profesionální tak amatérské sféře je trendem už několik
desítek let přechod od analogu k digitálním hlasovým módům. Ač mají analogové systémy stále svoje využití, praxe ukazuje, že digitální rádiové sítě
jsou spolehlivější, škálovatelnější a koncové radiostanice dokáží být operované i méně zkušeným radistou.

Myšlenka přenášet informaci rádiem v~digitální podobě je stejně stará jako rádio samo. Radiotelegrafie využívající Morseovu abecedu je v~podstatě
digitálním přenosem, neboť vysílač pouze zapíná a vypíná nosnou vlnu, a~přenáší tak dva diskrétní stavy. Hlasová komunikace se však po většinu
20.~století opírala výhradně o~analogové modulace AM a~FM. Digitální přenos dat se v~rádiovém provozu prosadil nejprve v~podobě dálnopisného provozu
RTTY a~v~80.~letech také paketového rádia založeného na protokolu AX.25. Přenos digitalizované řeči se stal prakticky použitelným až v~90.~letech,
kdy rozvoj číslicového zpracování signálu umožnil kompresi řeči do datových toků v~řádu jednotek kbit/s. V~této době vznikly první profesionální systémy
určené zejména pro bezpečnostní a~záchranné složky, jako jsou Tetrapol nebo TETRA. Po roce 2000 se objevily levnější úzkopásmové standardy DMR, dPMR
a~NXDN určené pro širší okruh komerčních uživatelů. Ve stejném období začali vlastní digitální módy využívat i~radioamatéři, nejprve D-STAR a~později
System Fusion, a~hojně přebírají také profesionální standard DMR. Většina uvedených systémů však spoléhá na proprietární hlasové kodeky rodiny AMBE, což vedlo
ke vzniku zcela otevřeného protokolu M17 s~kodekem Codec~2. Současným trendem v~profesionální sféře je přechod na širokopásmové sítě LTE a~5G se službami pro
kritickou komunikaci (MCPTT), které úzkopásmové systémy postupně doplňují. Přehled vybraných standardů uvádí tabulka~\ref{tab:digi-standardy}.

\begin{table}[H]
\centering
\small
\caption{Přehled vybraných standardů digitální hlasové rádiové komunikace}
\label{tab:digi-standardy}
\begin{tabular}{@{}l c l l l l@{}}
\toprule
\textbf{Standard} & \textbf{Rok} & \textbf{Přístup} & \textbf{Šířka kanálu} & \textbf{Modulace} & \textbf{Kodek} \\
\midrule
\multicolumn{6}{@{}l}{\textit{Profesionální systémy}} \\
\addlinespace[2pt]
Tetrapol        & 1992 & FDMA           & 12,5 kHz    & GMSK          & proprietární \\
TETRA           & 1995 & TDMA (4 sloty) & 25 kHz      & $\pi/4$-DQPSK & ACELP \\
APCO P25        & 1995 & FDMA           & 12,5 kHz    & C4FM          & IMBE \\
DMR             & 2005 & TDMA (2 sloty) & 12,5 kHz    & 4FSK          & AMBE+2 \\
dPMR            & 2006 & FDMA           & 6,25 kHz    & 4FSK          & AMBE+2 \\
NXDN            & 2006 & FDMA           & 6,25 kHz    & 4FSK          & AMBE+2 \\
LTE/5G MCPTT    & 2016 & OFDMA          & 1,4--20 MHz & QPSK/QAM      & AMR-WB \\
\midrule
\multicolumn{6}{@{}l}{\textit{Radioamatérské systémy}} \\
\addlinespace[2pt]
D-STAR          & 2001 & FDMA           & 6,25 kHz    & GMSK          & AMBE \\
System Fusion   & 2013 & FDMA           & 12,5 kHz    & C4FM          & AMBE+2 \\
M17             & 2019 & FDMA           & 9 kHz       & 4FSK          & Codec 2 \\
\bottomrule
\end{tabular}
\medskip
\begin{minipage}{0.9\linewidth}
\smallskip
\footnotesize
Pozn.: Rok udává vydání standardu, případně uvedení prvních zařízení či sítí do provozu.
Kodeky AMBE, AMBE+2 a~IMBE jsou proprietární kodeky, Codec~2 je otevřený.
Vokodér systému Tetrapol je typu RP-CELP, ve specifikaci se nazývá jednoduše
\uv{Tetrapol Vocoder}, protože byl specifikován pouze pro účely systému.
\end{minipage}
\end{table}

Velký užitek nachází digitální radiokomunikace například v koordinaci složek integrovaného záchranného systému, kde plní klíčovou roli pro jejich správné fungování.
V ČR vedle sebe funguje hned několik rádiových sítí používaných různými složkami IZS. Tou nejvýznamnější u nás je celostátní síť PEGAS postavená na technologii Tetrapol
a~provozovaná českou poštou pro potřeby ministerstva vnitra, která funguje jako páteřní síť IZS, ke které se ostatní sítě připojují. Je zapojená do všech operačních
středisek policie, hasičů a ZZS. Další celostátní sítí je Analogová rádio síť (ARS) HZS ČR. Ta slouží jako záložní síť hasičských jednotek k síti PEGAS.
Dále kromě pagerů a LTE, které jsou provozované soukromníky, jsou všechny sítě pouze na regionální úrovni.

\begin{table}[H]
\newcolumntype{P}[1]{>{\raggedright\arraybackslash}p{#1}}
\newcolumntype{L}{>{\raggedright\arraybackslash}X}
\centering
\small
\renewcommand{\arraystretch}{1.2}
\setlength{\tabcolsep}{5pt}
\newlength{\oblastw}\settowidth{\oblastw}{Jihomoravský}
\caption{Přehled vybraných rádiových sítí IZS}
\label{tab:site-izs}
\begin{tabularx}{\textwidth}{@{} P{\oblastw} L L L P{0.15\textwidth} @{}}
\toprule
\textbf{Oblast} & \textbf{Síť / technologie} & \textbf{Kdo ji používá} & \textbf{Detail} & \textbf{Stav} \\
\midrule
Celá ČR & \textbf{PEGAS}\newline(Tetrapol)
  & PČR, HZS, ZZS, AČR, ČEPS a~další menší složky IZS
  & Páteřní síť IZS
  & V~provozu\newline od r.~1995 \\
\addlinespace
Celá ČR & \textbf{ARS HZS ČR}\newline(analogová radiová síť)
  & HZS ČR a~jednotky SDH obcí a~podniků
  & Úzkopásmová FM se selektivní volbou ZVEI-1
  & V~provozu\newline od r.~2004, první pokusy od 50.~let \\
\addlinespace
Celá ČR & \textbf{Pagerové sítě}
  & SDH a~ZZS
  & Jednosměrný přenos čísel a~textu, svolávání jednotek
  & V~provozu \\
\addlinespace
Celá ČR & \textbf{Veřejné mobilní sítě}\newline(LTE/5G)
  & Všechny složky IZS
  & Datové služby, které úzkopásmový PEGAS nezvládá
  & V~provozu \\
\addlinespace
Praha & \textbf{Městský rádiový systém (MRS)}\newline(TETRA)
  & Městská policie, Dopravní podnik a~Technická správa komunikací
  & Součinnost účastníků v~rámci města
  & V~provozu\newline od r.~2000 \\
\addlinespace
Liberecký kraj & \textbf{DMR síť ZZS LK}\newline(DMR)
  & ZZS Libereckého kraje
  & Převaděče hlavně v~horském terénu, kde PEGAS nemá pokrytí
  & V~provozu \\
\addlinespace
Brno & \textbf{TETRA síť MP~Brno}\newline(TETRA)
  & Městská policie Brno
  & Běží vedle PEGASu a~DMR sítě ZZS JmK
  & V~provozu \\
\addlinespace
Jihomoravský kraj & \textbf{DMR síť ZZS JmK}\newline(DMR)
  & ZZS Jihomoravského kraje
  & DMR Tier~II, šifrování ARC4 (40~bitů)
  & V~provozu \\
\addlinespace
Břeclav, Znojmo, Blansko & \textbf{DMR sítě MP}\newline(DMR)
  & Městské policie Břeclav, Znojmo a~Blansko
  & DMR Tier~II, radiostanice s~GPS a~datovými službami
  & V~provozu \\
\bottomrule
\end{tabularx}
\end{table}

% --- DMR
\chapter{Standard DMR}

Systém DMR vznikl jako otevřený standard Evropského institutu pro telekomunikační normy (ETSI) určený především pro profesionální mobilní rádiové sítě.
Mezi radioamatéry se rozšířil díky dostupnosti levných ručních stanic a převaděčů a díky možnosti propojovat převaděče a uživatele po celém světě prostřednictvím internetu.
Právě propojení převaděčů s IP infrastrukturou otevírá prosto pro spolupráci s dalšími komunikačními systémy, včetně IP telefonie.

\section{Úrovně standardu (Tier I-III)}
Standard DMR rozlišuje celkem 3 úrovně, označované jako \textit{Tier I}, \textit{Tier II} a \textit{Tier III}.
Úrovně se liší zamýšleným nasazením, licencováním a složitostí síťové architektury.
Všechny však vycházejí ze společného rádiového rozhraní.


\textit{Tier I} je určen pro bezlicenční provoz v pásmu 446 MHz jako digitální ekvivalent analogových radiostanic PMR (Personal Mobile Radio).
Zařízení této úrovně pracují s nízkým vyzářeným výkonem (0.5W), omezeným počtem kanálů a striktně bez převaděčů.
Komunikace probíhá přímo mezi stanicemi v režimu point to point a provoz hovoru běží kontinuálně po jednom kanálu o šířce pásma 12.5 kHz. 


\textit{Tier II} je nejrozšířeněji používanou úrovní standardu a bude také předmětem této práce. Pokrývá nasazení v licencovaných pásmech VHF a UHF ($\sim$30 MHz - $\sim$1 GHz).
Využívá dvouslotový časový multiplex (TDMA) v kanálu o šířce 12.5 kHz, díky tomu lze na jednom kmitočtu vést dva nezávislé hovory.  Tier II podporuje provoz přes převaděče,
skupinová i individuální volání, krátké textové zprávy a paketová data.


\textit{Tier III} rozšiřuje Tier II o trunkový provoz. Rádiové kanály v systému nejsou pevně přiřazeny skupinám uživatelů, ale jsou přidělovány dynamicky podle potřeby.
Řízení zajišťuje vyhrazený řídicí kanál, přes který se stanice registrují do sítě a žádají o přidělení provozního kanálu.
Tento přístup umožňuje efektivně obsloužit velký počet uživatelů s omezeným počtem kmitočtů a podporuje rozsáhlé sítě s více základnovými stanicemi.

\section{Fyzická vrstva a přístup k médiu}
Pro přenos informace používá DMR čtyřstavovou frekvenční modulaci 4FSK, kde každý symbol nese 2 bity informace. Modulační rychlost je 4800 symbolů (Baudů) za vteřinu,
což činí teoretickou přenosovou rychlost 9600 bitů/s. Frekvenční skoky od středu kanálu jsou zaznamenány v následující tabulce:

\begin{table}[htbp]
\shorthandoff{-}
\centering
\caption{Mapování symbolů a bitů na odchylky modulace}
\label{tab:dibit}
\begin{tabular}{|c|c|c|c|}
\hline
\multicolumn{2}{|c|}{\textbf{Informace}} & \multirow{2}{*}{\textbf{Symbol}} & \multirow{2}{*}{\textbf{4FSK frekvenční odchylka [kHz]}} \\
\cline{1-2}
\textbf{Bit 1} & \textbf{Bit 0} & & \\
\hline
0 & 1 & $+3$ & $+1{,}944$ kHz \\
\hline
0 & 0 & $+1$ & $+0{,}648$ kHz \\
\hline
1 & 0 & $-1$ & $-0{,}648$ kHz \\
\hline
1 & 1 & $-3$ & $-1{,}944$ kHz \\
\hline
\end{tabular}
\end{table}

Šířka jednoho kanálu je 12,5 kHz, kde jeden kanál může obsloužit dva uživatele současně díky dvouslotovému časovému multiplexu TDMA
(Time Division Multiple Access). Střída (duty cycle) stanice je tak pouze 50\%, což má za následek vyšší energetickou efektivitu.

\section{Struktura rámců a signalizace}
\ldots

\section{Hlasové služby}
\ldots

\section{Adresace a typy volání}
\ldots

\section{Hlasový kodek AMBE+2}
\ldots

\section{Datové služby}
\ldots

% --- ASTERISK
\chapter{Telefonní sítě a Asterisk}
\ldots

% --- NÁVRH
\chapter{Návrh řešení spolupráce systému Asterisk s protokolem DMR}
\ldots

\section{Existující řešení propojení telefonie a rádia}
\ldots

\section{Síť BrandMeister}
\ldots

\subsection{Homebrew repeater protokol}
\ldots

\chapter{Regulatorní a etický rámec}
\label{chap:regule}
Tady chci napsat něco o ČTÚ a jeho regulacích na pásmech kde se DMR provozuje.
Dále etiku implementace takového mostu i přes výslovný zákaz autorů sítě BrandMeister o přemosťování komunikace mezi sítěmi.
Možná něco o uzavřenosti AMBE kodeku pokud se to nebude opakovat s tím co bude v kapitole o AMBE kodeku.

% --- IMPLEMENTACE
%\chapter{Implementace řešení}
%\ldots

% --- TESTOVÁNÍ
%\chapter{Testování a výsledky}
%\ldots

% --- ZÁVĚR
\chapter{Závěr}
\ldots

% --- Seznam použitých zdrojů ----------------------------------
% Kostra používá BibTeX (funguje out-of-the-box). Pro citace dle
% ČSN ISO 690 (což FEL/telekom obvykle chce) zvaž místo toho balíček
% biblatex se stylem iso-numeric nebo iso-authoryear.
\bibliographystyle{amsalpha}
\bibliography{literatura}        % zdroje v souboru literatura.bib

% --- Přílohy --------------------------------------------------
\appendix
% \printindex                    % rejstřík, používáš-li jej

% --- Vložení zadání z KOS jako kapitolu -----------------------
% \ctutemplate{specification.as.chapter}   % odkomentuj spolu se specification-file

\end{document}
